\section{Resolution of the recurrent difficulties}

This section revisits the main obstacles reported across the proof packets.  A
``resolution'' means either a complete generic lemma or a precise reduction to
an application-specific producer.  It does not mean that every deep theorem
has thereby been proved.

\subsection{Semantic algorithm versus object-language derivation}

\paragraph{Difficulty.}
A Lean function is shown to return a correct answer, but the target is a term of
\(\Prf\;\PRAAx\).  Re-running the function inside Lean does not create such a
term.

\paragraph{Resolution.}
Require every graph node to carry a local \(\KPrf\) or direct
\(\Prf\;\PRAAx\) derivation.  Compose native nodes with
Theorem~\ref{thm:dag-compile}; import direct nodes by
Theorem~\ref{thm:import}; export the terminal result by
Theorem~\ref{thm:export}.  Reference execution remains useful for generating
terms and detecting errors, but it is not the logical bridge.

\subsection{Substitution and de Bruijn bookkeeping}

\paragraph{Difficulty.}
Generic local lemmas must be specialized under binders, and hand-written
renaming equations dominate the Lean work.

\paragraph{Resolution.}
Target \(\KPrf\), where \path{KPrf.subst}, shifting, weakening, and cut have
already been proved.  Normalize arbitrary native formulas with \path{unQ}.
Keep application nodes at a small number of free variables and specialize them
before graph folding.  The remaining work is defining the correct substitution,
not re-proving its admissibility.

\subsection{A generic graph compiler}

\paragraph{Difficulty.}
Earlier frameworks define a proof graph but leave the compiler soundness theorem
or object-level fold unfinished.

\paragraph{Resolution.}
Theorem~\ref{thm:list-cut} and Theorem~\ref{thm:dag-compile} reduce the compiler
to checked \(\KPrf\) cut.  The graph is a metalevel finite list, not an encoded
object-level program.  A Lean implementation needs a dependent schedule type,
a membership router, and the direct list-cut proof; no new logical soundness
theorem is required.

\subsection{Bounded formulas inside quantifier-free induction}

\paragraph{Difficulty.}
An invariant is written with bounded universal or existential quantifiers and
therefore is not syntactically quantifier-free.

\paragraph{Resolution.}
Compile each fixed bounded formula to a Boolean fold and use the equation
\(\chi=1\) as the invariant.  Prove the Booleanization theorem once.  When a
witness is needed later, store it in the state rather than relying on an
existential inside the induction formula.

\subsection{Changing parameters and changing precision}

\paragraph{Difficulty.}
The statement to be proved at stage \(t\) depends on a precision or domain
parameter computed from later stages, so a direct induction has a changing
formula.

\paragraph{Resolution.}
Use the backward-demand/forward-replay construction.  The reverse trace computes
and codes the whole finite schedule.  The forward invariant refers to a lookup
in that fixed schedule.  Cyclic local requirements are closed first by
Theorem~\ref{thm:finite-closure}.

\subsection{Cyclic demand constraints}

\paragraph{Difficulty.}
Precision requirements depend mutually on one another, and naive iteration has
no termination proof.

\paragraph{Resolution.}
Work in the explicit finite box imposed by the theorem's resource envelope.
Iterate \(x\mapsto x\join F(x)\).  The rank sum strictly increases at each
nonfixed step, giving a primitive-recursive fuel and the least prefixed point.
If the transformer leaves the box, return a certified resource failure rather
than clipping.

\subsection{Noncanonical witnesses and singularities}

\paragraph{Difficulty.}
A proof asks for a root, eigenvector, chart, decomposition, or quotient
representative that cannot be selected continuously or uniformly at singular
inputs.

\paragraph{Resolution.}
Use finite relational clouds and backward viability.  Export only an invariant
observer, or an approximate observer with a diameter bound.  Introduce an exact
label only with a positive gap certificate.  At a collision, export cluster
information instead of a fictitious canonical branch.

\subsection{From finite cardinality to a search cutoff}

\paragraph{Difficulty.}
A theorem says that a solution set is finite or has at most \(m\) elements, and
a later algorithm treats this as a bound on the numerical codes of solutions.

\paragraph{Resolution and limitation.}
A cardinality bound alone is not a code-height bound.  The framework proves the
following exact equivalence at the right level.

\begin{theorem}[Height--list conversion]\label{thm:height-list}
Let \(P(x)\) be decidable.
\begin{enumerate}
\item If \(H\) satisfies \(P(x)\to x\le H\), then bounded enumeration through
\(H\) produces a complete finite list of all solutions.
\item If a coded finite list \(L\) is sound and complete for \(P\), then
\(H=\max(L)\) is a code-height bound, with \(H=0\) for the empty list.
\end{enumerate}
\end{theorem}

\begin{proof}
The first part is Theorem~\ref{thm:enumeration}.  For the second, completeness
gives membership of every solution in \(L\); a bounded fold proves every list
entry is at most its maximum.
\end{proof}

A bare statement ``there are at most \(m\) solutions'' supplies neither object.
General Faltings or Siegel arguments therefore remain conditional until a
suitable effective height theorem or complete-list producer is provided.

\subsection{Ineffective constants}

\paragraph{Difficulty.}
An analytic or number-theoretic theorem uses ``for sufficiently large \(x\)''
with a constant whose proof is ineffective or whose value is not extracted.

\paragraph{Resolution and limitation.}
Separate the theorem into:

\begin{enumerate}
\item a uniform verifier/finalizer theorem parameterized by a proposed threshold
\(X_0\), constants, and finite proof certificates;
\item a producer theorem asserting that suitable data exist and are supplied.
\end{enumerate}

The first can often be proved in iPRA.  The second is not solved by calling its
output a certificate.  For finite exceptional sets with explicit budgets,
Theorem~\ref{thm:common-window} gives a common search window.  Without explicit
budgets or thresholds, the framework records an open producer obligation.

\subsection{Classical compactness and local-to-global arguments}

\paragraph{Difficulty.}
A proof uses compactness to choose a finite subcover, a Lebesgue number, or a
uniform modulus, but the formal statement provides no finite data.

\paragraph{Resolution.}
Use one of four interfaces: supplied finite cover, bounded cover enumerator,
finite subcover certificate, or explicit mesh/Lebesgue-number bound.  Compile
local proofs as clause nodes and overlaps as compatibility nodes.  A nerve or
partition data structure is finite.  The framework does not claim that every
classically compact presentation has a PR cover producer.

\subsection{Completed limits and infinite paths}

\paragraph{Difficulty.}
A classical proof forms an infinite sequence of choices or passes to a limit.

\paragraph{Resolution.}
State the theorem at arbitrary finite precision or horizon.  Use finite
viability for paths and a motion-tail bound for Cauchy coherence.  External
adequacy reconstructs the ideal limit.  When no modulus is available, use an
explicit metastable formulation.  Do not assert a uniform Cauchy modulus merely
because a classical limit exists.

\subsection{Quotients and canonical forms}

\paragraph{Difficulty.}
The development tries to normalize every quotient representative, leading to a
large and sometimes impossible canonicalization problem.

\paragraph{Resolution.}
Use an equivalence relation and stable observers.  Normalize only where a
verified PR normalizer is useful.  The clause graph can transport an observer
through equivalence without choosing a representative.  This is particularly
effective for rational fractions, polynomial roots up to permutation, bases of
modules, and isomorphism classes of finite structures.

\subsection{Proof-size explosion}

\paragraph{Difficulty.}
Expanding a finite computation into individual logical steps produces enormous
proof terms.

\paragraph{Resolution.}
Prove generic trace, fold, and algebra lemmas once.  Represent long computations
by PR functions plus their defining equations and one induction.  Use symbolic
ring derivations for variable identities and certificate checking for concrete
large objects.  The graph compiler reuses local lemmas by substitution instead
of duplicating their proofs.

\subsection{What remains genuinely open}

The framework does not currently supply:

\begin{itemize}
\item a kernel-checked generic graph compiler implementation in the repository;
\item a complete checked library for pair/list coding, bounded Booleanization,
finite tables, rational order, and analytic error arithmetic in
\(\Prf\;\PRAAx\);
\item general effective Faltings or Siegel bounds;
\item uniform primitive-recursive Mordell--Weil representatives for arbitrary
input curves;
\item effective constants for every classical prime-distribution theorem;
\item exact selectors through arbitrary root or spectral collisions;
\item a bridge from other PRA formalizations in the repository to
\path{ConciseIPRA}'s \(\Prf\;\PRAAx\).
\end{itemize}

These are now localized obligations.  They no longer block the general finite
control and composition architecture.

\subsection{Status table}

\begin{longtable}{L{0.29\textwidth}L{0.18\textwidth}L{0.44\textwidth}}
\toprule
Issue & Framework status & Exact endpoint \\
\midrule
Graph composition & Resolved mathematically & List cut plus topological induction.\\
Parameter specialization & Checked upstream & \path{KPrf.subst}.\\
Native-to-iPRA export & Checked upstream plus one-line composition & \path{KPrf.toNPrf}; \path{native_iff_pra}.\\
Direct algebra import & Checked upstream plus one-line composition & \path{pra_iff_native}; \path{kprf_iff_nprf}.\\
Bounded search/enumeration & Resolved mathematically & Clocked trace with Boolean invariant.\\
Cyclic demand & Resolved mathematically & Finite monotone closure theorem.\\
Finite relational choice & Resolved mathematically & Backward viability and least path.\\
Stable observables & Resolved mathematically & Diameter and motion-tail theorems.\\
Exact branch labels & Conditional & Positive separation gap required.\\
Finite-set synchronization & Resolved with budgets & Common good finite window.\\
Finiteness to enumeration & Conditional & Code height or complete list required.\\
General deep Diophantine producers & Open/application-specific & Explicit producer interface.\\
Ineffective analytic constants & Open/application-specific & Explicit threshold/certificate interface.\\
New Lean compiler & Outlined, not checked here & Dependency plan in Section~\ref{sec:lean-plan}.\\
\bottomrule
\end{longtable}
